% Stany pracy falownika polmostkowego - drogi przeplywu pradu obciazenia.
% Wstawiany wewnatrz srodowiska figure; podpis i etykieta w pliku rozdzialu.
% #1 - zamknieta droga pradu, podana zgodnie z rzeczywistym kierunkiem pradu
% #2 - nazwa przewodzacego elementu (wyrozniona kolorem)
% #3 - strzalki kierunku pradu: lista x/y/kat (kat w stopniach)
\providecommand{\polmostekstan}[3]{%
\begin{circuitikz}[european, scale=0.8, transform shape, font=\normalsize,
    sciezka/.style={draw=red!25, line width=7pt, line join=round}]
  \useasboundingbox (-1.3,-0.35) rectangle (7.0,6.35);
  % droga pradu (pod schematem)
  \path[sciezka] #1;
  % schemat polmostka
  \draw (0,0) to[V, l=$U_\mathrm{DC}$, invert] (0,6);
  \draw (0,6) -- (6.2,6);
  \draw (0,0) -- (6.2,0);
  \node[nmos, anchor=D] (Q1) at (2.2,6) {};
  \node[nmos, anchor=S] (Q2) at (2.2,0) {};
  \draw (Q1.S) -- (Q2.D);
  \node[left] at (Q1.G) {$Q_1$};
  \node[left] at (Q2.G) {$Q_2$};
  \draw (3.1,3) to[D*, l_=$D_1$, *-*] (3.1,6);
  \draw (3.1,0) to[D*, l_=$D_2$, *-*] (3.1,3);
  \draw (2.2,3) node[circ]{} -- (3.1,3);
  \node[below left] at (2.2,3) {a};
  \draw (6.2,6) to[C, l=$C_1$, *-] (6.2,3) to[C, l=$C_2$, -*] (6.2,0);
  \draw (6.2,3) node[circ]{} node[right]{0};
  \draw (3.1,3) to[generic, l=$Z_\mathrm{o}$] (6.2,3);
  % element przewodzacy
  \node[draw=red!80!black, very thick, rounded corners=2pt, fill=white,
        text=red!80!black, inner sep=2pt, font=\normalsize\bfseries]
        at (4.65,5.2) {#2};
  % strzalki kierunku pradu (nad schematem)
  \foreach \x/\y/\a in {#3}
    \draw[red!80!black, line width=1.3pt, -{Latex[length=3.4mm, width=2.8mm]}]
      ([shift={(\a:-0.32)}]\x,\y) -- ([shift={(\a:0.32)}]\x,\y);
\end{circuitikz}}%
\begin{subfigure}[b]{0.48\textwidth}
  \centering
  \polmostekstan{(2.2,6) -- (2.2,3) -- (6.2,3) -- (6.2,6) -- cycle}{$Q_1$}{2.2/3.6/-90, 3.5/3/0, 6.2/3.6/90, 4.2/6/180}
  \caption{$i_\mathrm{o} > 0$, przewodzi $Q_1$,
    $u_\mathrm{o} = +U_\mathrm{DC}/2$}
\end{subfigure}\hfill
\begin{subfigure}[b]{0.48\textwidth}
  \centering
  \polmostekstan{(3.1,0) -- (3.1,3) -- (6.2,3) -- (6.2,0) -- cycle}{$D_2$}{3.1/2.55/90, 3.5/3/0, 6.2/2.45/-90, 4.6/0/180}
  \caption{$i_\mathrm{o} > 0$, przewodzi $D_2$,
    $u_\mathrm{o} = -U_\mathrm{DC}/2$}
\end{subfigure}

\vspace{5mm}
\begin{subfigure}[b]{0.48\textwidth}
  \centering
  \polmostekstan{(6.2,3) -- (2.2,3) -- (2.2,0) -- (6.2,0) -- cycle}{$Q_2$}{5.8/3/180, 2.2/2.4/-90, 4.2/0/0, 6.2/0.6/90}
  \caption{$i_\mathrm{o} < 0$, przewodzi $Q_2$,
    $u_\mathrm{o} = -U_\mathrm{DC}/2$}
\end{subfigure}\hfill
\begin{subfigure}[b]{0.48\textwidth}
  \centering
  \polmostekstan{(6.2,3) -- (3.1,3) -- (3.1,6) -- (6.2,6) -- cycle}{$D_1$}{5.8/3/180, 3.1/3.45/90, 4.6/6/0, 6.2/5.45/-90}
  \caption{$i_\mathrm{o} < 0$, przewodzi $D_1$,
    $u_\mathrm{o} = +U_\mathrm{DC}/2$}
\end{subfigure}
